\theory{Morse theory}{How can the critical points of an ordinary height function reveal the topology of a manifold?}{Morse theory studies smooth functions whose critical points are nondegenerate. As a level rises, the manifold changes only when it passes a critical point, and each critical point attaches a cell or handle.}{Differential topology, homology, geodesics, dynamical systems, optimization, and mathematical physics.}{Critical points, Hessian indices, and gradient flows convert local differential information into a global cell or handle decomposition.}

\paragraph{History and the landscape picture}
Imagine a mountain landscape. At most heights, contour lines change smoothly. At a valley a new component appears; at a saddle two contours join or split; at a peak a component disappears. Marston Morse made this picture precise. Earlier topographical ideas came from Cayley and Maxwell, and Morse first applied the method to geodesics and energy functionals \cite{morse_theory_ref1}.

A smooth function $f:M\to\mathbb R$ has a critical point where its derivative is zero. It is nondegenerate when the Hessian has no zero eigenvalue. The number of negative Hessian directions is the Morse index. A Morse function has only nondegenerate critical points, and a small generic perturbation usually makes a function Morse.

\end{historylist}
\section*{3. Theory}
\subsection*{Core theory}
\paragraph{The Morse lemma}
Near a critical point of index $\lambda$, suitable coordinates make
\[
f=f(p)-x_1^2-\cdots-x_\lambda^2+x_{\lambda+1}^2+\cdots+x_n^2.
\]
Thus every nondegenerate critical point looks locally like a standard quadratic saddle. Below the critical value the sublevel set $M_a=\{x:f(x)\leq a\}$ changes only by attaching a $\lambda$-dimensional cell when the level crosses the critical point.

\paragraph{Topology from critical points}
Climbing the function builds a CW structure on $M$. One cell of dimension $\lambda$ appears for each critical point of index $\lambda$. Consequently the number $C_\lambda$ of such critical points is at least the Betti number $b_\lambda$. The alternating sum gives
\[
\chi(M)=\sum_\lambda(-1)^\lambda C_\lambda.
\]
Morse inequalities strengthen this to bounds on alternating sums of critical-point counts. A sphere can be built from one minimum and one maximum; a torus needs a minimum, two saddles, and a maximum.

\paragraph{Closed surfaces and handles}
On a compact surface, the index-zero critical points add disks, index-one points add bands, and index-two points cap off the surface. Counting them distinguishes the sphere, torus, and higher-genus surfaces. Handle decompositions in higher dimensions are the manifold version of the same construction.

Morse homology forms a chain complex generated by critical points, with differential counting gradient-flow lines between points whose indices differ by one. Its homology is isomorphic to singular homology, so analysis of a gradient field computes topological holes.

\paragraph{Morse--Bott theory and other applications}
Morse--Bott theory permits the critical set to be a smooth submanifold rather than isolated points, while requiring a nondegenerate Hessian in normal directions. This is useful when symmetry forces critical points to occur in families.

Morse theory analyzes geodesics as critical points of the energy on path spaces and was used in Bott's proof of periodicity. In optimization, minima, maxima, and saddles describe landscapes; in physics, energy functions and instantons produce similar cell attachments. The complex analogue is Picard--Lefschetz theory \cite{morse_theory_ref2}.

\paragraph{What the theory depends on and where it is useful}
Morse theory depends on calculus, differential topology, manifolds, flows, and homology. It helps handle and cobordism theory construct manifolds, Hodge theory study energy minimizers, and dynamics understand gradient trajectories.

The method requires smoothness and nondegenerate critical points. Critical-point counts give bounds and decompositions, but two different functions on the same manifold may produce very different numbers of critical points.

\section*{4. Important theorems and results}
\begin{enumerate}[leftmargin=*,itemsep=0.35em]
\item \textbf{Morse lemma.} Every nondegenerate critical point has a standard quadratic normal form.

\item \textbf{Morse inequalities.} Critical points of each index bound the corresponding Betti numbers and satisfy alternating inequalities.

\item \textbf{Handle decomposition theorem.} A compact smooth manifold admits a handle decomposition associated with a Morse function.

\item \textbf{Morse homology theorem.} The homology of the Morse complex agrees with singular homology.

\item \textbf{Morse--Bott generalization.} Nondegenerate critical submanifolds replace isolated critical points \cite{morse_theory_ref3}.

Together these results convert a smooth function into a topological construction. The Morse lemma gives local quadratic coordinates, the inequalities compare critical points with holes, and handle decomposition explains how the manifold is assembled. Morse homology is important because it produces the same invariant as singular homology from gradient trajectories, while the Morse--Bott version shows how the method survives when critical points occur in families.
\end{enumerate}

\theoryapplications
\theoryconnections{morse theory}{%
\item Differential calculus supplies critical points and Hessians, differential topology supplies smooth manifolds, and algebraic topology supplies homology and Betti numbers.
\item Gradient flows and handle decompositions connect local derivatives with global shape.
}{%
\item Morse theory helps compute the topology of manifolds from critical points and supports optimization, geometry, and dynamical systems.
\item Its inequalities give reliable information even when the full manifold or all of its critical points is difficult to describe.
}
\section*{7. Sample Questions}
\emph{Exercise reference:} \cite{morse_theory_ref1}. The cited edition is the source for the exercise set; consult its Exercises section for the official statement and numbering.
\begin{enumerate}[leftmargin=*,itemsep=0.9em]
\item \textbf{Exercise 1.} What is the Morse index of a local minimum?

\emph{Solution.} Zero, because the Hessian is positive definite and has no negative direction.

\item \textbf{Exercise 2.} What does an index-one critical point do to a surface?

\emph{Solution.} It attaches a one-dimensional handle or band. In a landscape it is a saddle where contour components join or split.

\item \textbf{Exercise 3.} How many critical points can a height function on a sphere have in the simplest Morse picture?

\emph{Solution.} One minimum and one maximum, of indices zero and two. Their alternating sum is $1+1=2$, the Euler characteristic of the sphere.

\item \textbf{Exercise 4.} Why do Morse inequalities imply topological information?

\emph{Solution.} Homology ranks count holes, and each homology generator requires enough cells. Critical points create those cells, so their number bounds the holes.

\item \textbf{Exercise 5.} Why is Morse--Bott theory needed?

\emph{Solution.} Symmetry may make an entire circle or submanifold critical. Morse--Bott theory handles the family while retaining nondegeneracy transverse to it.

\end{enumerate}
\section*{8. References}
\addcontentsline{toc}{section}{References}

\begin{thebibliography}{9}
\bibitem{morse_theory_ref1} John Milnor, \emph{Morse Theory}, Princeton University Press, 1963.
\bibitem{morse_theory_ref2} Raoul Bott, \emph{Morse Theory Indomitable}, Publications of the Mathematical Society of Japan, 1988.
\bibitem{morse_theory_ref3} Matthias Schwarz, \emph{Morse Homology}, Birkhauser, 1993.
\end{thebibliography}
